\documentclass[../../../include/open-logic-section]{subfiles}

\begin{document}

\olfileid[hi]{sth}{spine}{recursion}
\olsection{परिमितेतर पुनरावर्तन प्रमेय}

सबसे पहले हमें यह पुष्टि करनी होगी कि \olref[valpha]{defValphas} एक सफल
परिभाषा है। अधिक सामान्यतः हमें सिद्ध करना है कि (परिमितेतर) पुनरावर्तन
द्वारा परिमितेतर परिभाषा देने का प्रत्येक प्रयास सफल होगा। यही इस अनुभाग का
उद्देश्य है।

\emph{चेतावनी: यह कठिन सामग्री है।} फिर भी व्यापक शिक्षा काफ़ी सरल है:
परिमितेतर आगमन और प्रतिस्थापन मिलकर परिमितेतर पुनरावर्तन के (कई रूपों की)
वैधता की गारंटी देते हैं।\footnote{स्मरण रहे: सभी सूत्रों और पदों में प्राचल
हो सकते हैं (जब तक स्पष्टतः अन्यथा न कहा गया हो)।}

\begin{defn}
	$\tau(x)$ को पद, $f$ को फलन और $\alpha$ को क्रमसूचक मानें। हम $f$ को
	$\tau$ का \emph{$\alpha$-सन्निकटन} तभी और केवल तभी कहते हैं जब
	$\dom{f} = \alpha$ और
	$(\forall \beta \in \alpha)f(\beta) =
	\tau(\funrestrictionto{f}{\beta})$, दोनों हों।
\end{defn}

\begin{lem}[परिबद्ध पुनरावर्तन]\ollabel{transrecursionfun}
किसी भी पद $\tau(x)$ और किसी भी क्रमसूचक $\alpha$ के लिए $\tau$ का एक
अद्वितीय $\alpha$-सन्निकटन है।
\end{lem}
\begin{proof}
हम दिखाएँगे कि किसी भी $\gamma \leq \alpha$ के लिए एक अद्वितीय
$\gamma$-सन्निकटन है।

पहले अद्वितीयता स्थापित करें। $g$ और $h$ को क्रमशः $\gamma$- और
$\delta$-सन्निकटन मानें। उनके निवेशों पर परिमितेतर आगमन दिखाता है कि किसी भी
$\beta \in \dom{g} \cap \dom{h} = \gamma \cap \delta =
\min(\gamma, \delta)$ के लिए $g(\beta) = h(\beta)$। अतः हमारे सन्निकटन
(यदि उनका अस्तित्व है) अद्वितीय हैं और सभी मानों पर सहमत हैं।

अस्तित्व स्थापित करने के लिए अब हम क्रमसूचकों $\delta \leq \alpha$ पर सरल
परिमितेतर आगमन (\olref[ordinals][opps]{simpletransrecursion}) का उपयोग करते
हैं।

रिक्त फलन तुच्छतः $\emptyset$-सन्निकटन है।

यदि $g$ कोई $\gamma$-सन्निकटन है, तो
$g \cup \{\tuple{\gamma, \tau(g)}\}$ कोई
$\ordsucc{\gamma}$-सन्निकटन है।

यदि $\gamma$ सीमा क्रमसूचक है और सभी $\delta < \gamma$ के लिए $g_\delta$
कोई $\delta$-सन्निकटन है, तो $g = \bigcup_{\delta \in \gamma} g_\delta$
रखें। यह फलन है, क्योंकि हमारे विभिन्न $g_\delta$ सभी मानों पर सहमत हैं।
और यदि $\delta \in \gamma$, तो
$g(\delta) = g_{\ordsucc{\delta}}(\delta) =
\tau(\funrestrictionto{g_{\ordsucc{\delta}}}{\delta}) =
\tau(\funrestrictionto{g}{\delta})$।

इससे परिमितेतर आगमन द्वारा प्रमाण पूर्ण होता है।
\end{proof}

यदि हम किसी फलन के बजाय \emph{पद} को परिभाषित करने की अनुमति लें, तो पिछले
परिणाम से परिबंध $\alpha$ हटाया जा सकता है। निम्न परिणाम के कथन और प्रमाण
में, जब $\sigma$ पद हो, हम
$\funrestrictionto{\sigma}{\alpha} = \Setabs{\tuple{\beta,
\sigma(\beta)}}{\beta \in \alpha}$ रखते हैं।

\begin{thm}[सामान्य पुनरावर्तन]\ollabel{transrecursionschema}
किसी भी पद $\tau(x)$ के लिए हम स्पष्टतः ऐसा पद $\sigma(x)$ परिभाषित कर सकते
हैं कि किसी भी क्रमसूचक~$\alpha$ के लिए
$\sigma(\alpha) = \tau(\funrestrictionto{\sigma}{\alpha})$।
\end{thm}

\begin{proof}
प्रत्येक $\alpha$ के लिए \olref{transrecursionfun} से~$\tau$ का अद्वितीय
$\alpha$-सन्निकटन $f_\alpha$ है। $\sigma(\alpha)$ को
$f_{\ordsucc{\alpha}}(\alpha)$ के रूप में परिभाषित करें। अब:
	\begin{align*}
		\sigma(\alpha) &= 
		f_{\ordsucc{\alpha}}(\alpha) \\&= 
		\tau(\funrestrictionto{f_{\ordsucc{\alpha}}}{\alpha}) \\&= 
		\tau(\Setabs{\tuple{\beta, f_{\ordsucc{\alpha}}(\beta)}}{\beta \in \alpha}) \\&= 
		\tau(\Setabs{\tuple{\beta, f_{\ordsucc{\beta}}(\beta)}}{\beta \in \alpha})\\&=
		\tau(\funrestrictionto{\sigma}{\alpha})
	\end{align*}
यह ध्यान रखते हुए कि \olref{transrecursionfun} की तरह सभी
$\beta < \alpha$ के लिए
$f_{\ordsucc{\beta}}(\beta) = f_{\ordsucc{\alpha}}(\beta)$।
\end{proof}
\noindent 
ध्यान दें कि \olref{transrecursionschema} एक \emph{योजना} है। निर्णायक रूप
से हम $\sigma$ से किसी फलन, अर्थात किसी विशेष प्रकार के \emph{समुच्चय}, को
परिभाषित करने की अपेक्षा नहीं कर सकते; क्योंकि तब $\dom{\sigma}$ समस्त
क्रमसूचकों का समुच्चय होता, जो बुराली--फ़ोर्टी विरोधाभास
(\olref[ordinals][basic]{buraliforti}) के विरुद्ध है।

फिर भी यह दिखाना शेष है कि \olref{transrecursionschema}, $V_\alpha$ की
हमारी परिभाषा का औचित्य सिद्ध करता है। यह तुरंत स्पष्ट न हो, पर परिमितेतर
पुनरावर्तन के एक अंतिम सरल रूप से स्पष्ट हो जाएगा।

\begin{thm}[सरल पुनरावर्तन]\ollabel{simplerecursionschema}
किन्हीं पदों $\tau(x)$ और $\theta(x)$ तथा किसी भी समुच्चय $A$ के लिए हम
स्पष्टतः ऐसा पद $\sigma(x)$ परिभाषित कर सकते हैं कि:
\begin{align*}
	\sigma(\emptyset) &= A\\
	\sigma(\ordsucc{\alpha}) &= \tau(\sigma(\alpha)) &&
		\text{किसी भी क्रमसूचक }\alpha\text{ के लिए}\\
	\sigma(\alpha) &= \theta(\ran{\funrestrictionto{\sigma}{\alpha}})&&
	\text{जब }\alpha\text{ सीमा क्रमसूचक हो}
\end{align*}
\end{thm}

\begin{proof}
हम पद $\xi(x)$ को इस प्रकार परिभाषित करने से आरंभ करते हैं:
%t $\xi(x) = A$ if $x$ is not a function whose domain is an ordinal;
%otherwise let:
\[
	\xi(x) = 
	\begin{cases}
			A &\text{यदि $x$ ऐसा फलन नहीं जिसका}\\
			&\hspace{1em}\text{प्रांत क्रमसूचक हो; अन्यथा:}\\
			\tau(x(\alpha)) & \text{यदि $\dom{x} = \ordsucc{\alpha}$}\\
			\theta(\ran{x}) & \text{यदि $\dom{x}$ सीमा क्रमसूचक हो}
	\end{cases}
\]
\olref{transrecursionschema} से ऐसा पद $\sigma(x)$ है कि प्रत्येक क्रमसूचक
$\alpha$ के लिए
$\sigma(\alpha) = \xi(\funrestrictionto{\sigma}{\alpha})$; इसके अतिरिक्त
$\funrestrictionto{\sigma}{\alpha}$ प्रांत $\alpha$ वाला फलन है। सरल
परिमितेतर आगमन (\olref[ordinals][opps]{simpletransrecursion}) से हम दिखाते
हैं कि $\sigma$ में अपेक्षित गुण हैं।

पहले, $\sigma(\emptyset) = \xi(\emptyset) = A$।

अगला,
$\sigma(\ordsucc{\alpha}) = \xi(\funrestrictionto{\sigma}{\ordsucc{\alpha}})
= \tau(\funrestrictionto{\sigma}{\ordsucc{\alpha}}(\alpha))
= \tau(\sigma(\alpha))$।

अंत में, जब $\alpha$ सीमा हो, तब
$\sigma(\alpha) = \xi(\funrestrictionto{\sigma}{\alpha})
= \theta(\ran{\funrestrictionto{\sigma}{\alpha}})$।
\end{proof}
\noindent
अब \olref[valpha]{defValphas} का औचित्य सिद्ध करने के लिए बस
$A = \emptyset$, $\tau(x) = \Pow{x}$ और $\theta(x) = \bigcup x$ लें।
अंततः इससे $V_\alpha$ की परिभाषा का औचित्य सिद्ध होता है!{}


\end{document}
